\documentclass[12pt]{article}
\def\activityid{F08-U-v2}\usepackage[letterpaper,margin=.65in,top=.60in,bottom=.65in]{geometry}
\usepackage[T1]{fontenc}
\usepackage[default]{sourcesanspro}
\usepackage{microtype,tikz,array,tabularx,fancyhdr,amsmath,amssymb}
\usepackage[hidelinks]{hyperref}
\usetikzlibrary{calc,arrows.meta}
\setlength{\parindent}{0pt}\setlength{\parskip}{7pt}
\setlength{\headheight}{0pt}\setlength{\footskip}{26pt}\setlength{\emergencystretch}{2em}
\pagestyle{fancy}\fancyhf{}\renewcommand{\headrulewidth}{0pt}\renewcommand{\footrulewidth}{.4pt}
\fancyfoot[L]{\scriptsize Bellingham Math Circle / Week 8 / \activityid}
\fancyfoot[R]{\scriptsize \thepage}
\definecolor{ink}{gray}{.12}\definecolor{soft}{gray}{.95}\color{ink}
\newcommand{\PageTitle}[2]{{\fontsize{10}{12}\selectfont\bfseries WEEK 8\quad STRATEGY LAB\hfill #1\par}\vspace{3pt}{\fontsize{26}{29}\selectfont\bfseries #2\par}\vspace{2pt}}
\newcommand{\NameLine}{{\small Name: \rule{2.65in}{.4pt}\hfill Date: \rule{1.35in}{.4pt}\par}\vspace{4pt}}
\newcommand{\Task}[2]{\par\vspace{3pt}{\large\bfseries #1}\par #2\par}
\newcommand{\Subhead}[1]{\par\vspace{3pt}{\large\bfseries #1}\par}
\newcommand{\RuleBox}[1]{\par\begingroup\setlength{\fboxsep}{8pt}\colorbox{soft}{\parbox{\dimexpr\linewidth-16pt\relax}{#1}}\endgroup\par}
\newcommand{\WriteBox}[2]{\par\textbf{#1}\par\vspace{2pt}\begin{tikzpicture}\draw[black!40,line width=.5pt] (0,0) rectangle (\dimexpr\linewidth-.5pt\relax,#2);\end{tikzpicture}\par}
\newcommand{\StudentFooter}{\vfill{\small Try it with objects. Explain by pointing, talking, or drawing.}\par}
\newcommand{\RookBoard}[3]{\begin{tikzpicture}[x=#3,y=#3]\draw[step=1,black!45] (0,0) grid (#1,#2);\draw[thick] (0,0) rectangle (#1,#2);\node[font=\Large] at ({#1-.5},.5){$\star$};\draw[-{Latex},thick] (0,-.35)--(#1,-.35);\draw[-{Latex},thick] ({#1+.35},#2)--({#1+.35},0);\end{tikzpicture}}
\newcommand{\Table}[3]{\begin{tikzpicture}[x=#3,y=#3]\draw[step=1,black!25] (0,0) grid (#1,#2);\draw[thick] (0,0) rectangle (#1,#2);\fill (0,0)circle(3pt);\draw[-{Latex},very thick] (0,0)--(.7,.7);\node[below] at ({#1/2},-.1){width #1};\node[rotate=90,above] at (-.2,{#2/2}){height #2};\end{tikzpicture}}
\newcommand{\GraphNodes}{\coordinate(A)at(0,0);\coordinate(B)at(3,0);\coordinate(C)at(1.5,2.6);\coordinate(D)at(1.5,.95);}
\newcommand{\Kfour}[1]{\begin{tikzpicture}[scale=#1]\GraphNodes\draw[very thick](A)--(B)--(C)--(A)--(D)--(B) (D)--(C);\foreach \n in {A,B,C,D}{\node[circle,draw,fill=white,minimum size=22pt,font=\bfseries]at(\n){\n};}\end{tikzpicture}}
\newcommand{\House}[1]{\begin{tikzpicture}[scale=#1]\coordinate(A)at(0,0);\coordinate(B)at(3,0);\coordinate(C)at(3,2);\coordinate(D)at(0,2);\coordinate(E)at(1.5,3.3);\draw[very thick](A)--(B)--(C)--(D)--(A) (D)--(E)--(C);\foreach \n in {A,B,C,D,E}{\node[circle,draw,fill=white,minimum size=22pt,font=\bfseries]at(\n){\n};}\end{tikzpicture}}

\begin{document}

\PageTitle{GRADES 4--5}{Three piles break the mirror}\NameLine
\RuleBox{Build three piles. On a turn, remove any positive number of counters from exactly ONE pile. Take the last counter to win. Use ordinary counters, not bundles, for actual play.}
\Task{1. Challenge an adult.}{Try piles \((1,2,3)\), then \((1,2,4)\). Change who moves first. Which starts seem to be traps for the player receiving them?}
\WriteBox{Record moves, replies, and a conjecture.}{1.05in}
\Task{2. Compare two kinds of balance.}{With only two piles, matching sizes lets you copy your opponent. With three piles, investigate \((1,1,1)\), \((1,1,2)\), \((1,2,3)\), and \((2,4,6)\). Is an even \emph{total} enough to make a trap?}
\WriteBox{Find evidence against a rule that almost works.}{1.0in}
\Task{3. Trade for bundles.}{Use labels 1, 2, and 4. Write each pile as a sum using each label at most once: for example, 5 is 4+1. For each label, count how many piles use it. Which starts have an even count for every bundle?}
\textbf{Try before turning the page.} An adult can make the bundle labels with scraps of paper.
\StudentFooter

\newpage

\PageTitle{GRADES 4--5}{Balance the binary columns}\NameLine
\RuleBox{Represent each pile using bundles 4, 2, 1. Write 1 if that bundle is used and 0 if not. Call a position \textbf{balanced} when every column contains an even number of 1s. Zero 1s counts as even.}
\begin{center}\renewcommand{\arraystretch}{1.5}\begin{tabular}{c|ccc@{\qquad}c|ccc}Pile&4&2&1&Pile&4&2&1\\\hline1&0&0&1&1&0&0&1\\2&0&1&0&2&0&1&0\\3&0&1&1&4&1&0&0\end{tabular}\end{center}
\Task{4. Repair the unbalanced position.}{Change just one pile in \((1,2,4)\) to a \emph{smaller} number, making every column even. What move does this mean with counters?}
\Task{5. Test your rule.}{For each start, either give a move to balance or say it is already balanced. Then play against someone trying to break your prediction.}
\begin{center}\renewcommand{\arraystretch}{1.8}\begin{tabular}{c|p{2.0in}|p{2.2in}}Start&4/2/1 records&Move or already balanced?\\\hline(2,4,6)&&\\(3,4,5)&&\\(1,5,7)&&\\(3,5,6)&&\end{tabular}\end{center}
\WriteBox{Can a legal move from balance leave every column even? Why?}{.85in}
\textbf{Vocabulary after discovery:} The column parities form the \emph{Nim sum}, also called bitwise XOR. A balanced position has Nim sum zero.
\StudentFooter

\newpage

\PageTitle{GRADES 4--5}{Why must the strategy work?}\NameLine
\Task{6. The highest wrong column.}{Suppose some column is odd. Find the largest bundle whose column is odd. At least one pile uses it. In that pile:}
\begin{enumerate}\item Keep all larger bundle choices unchanged.\item Remove this bundle.\item Choose the smaller bundles to make \emph{every} column even.\end{enumerate}
Try this recipe on \((3,4,5)\). Then try \((7,10,12)\), adding an 8-column.
\WriteBox{My moves, and why they fix all the columns:}{1.05in}
\Task{7. Two missing pieces of a proof.}{Why is the new pile \emph{smaller}, even if you add some smaller bundles? Why can the game not continue forever?}
\WriteBox{A reason that works for any pile sizes:}{1.05in}
\Task{8. Design a trap.}{Choose any first two piles. Can you choose the third to make balance? Could you do it in two different ways? Give your puzzle and its explanation to an adult.}
\WriteBox{My designed start and why its third pile is forced:}{.75in}
\StudentFooter

\end{document}
